\documentclass[12pt, a4paper]{article}

\usepackage[utf8]{inputenc}
\usepackage[T1]{fontenc}
\usepackage{geometry}
\usepackage{graphicx}
\usepackage{booktabs}
\usepackage{hyperref}
\usepackage{amsmath}
\usepackage{float}
\usepackage{enumitem}
\usepackage{caption}
\usepackage{listings}
\usepackage[table]{xcolor}
\usepackage{fancyhdr}
\usepackage{setspace}

\geometry{left=2.5cm, right=2.5cm, top=2.5cm, bottom=2.5cm}

\pagestyle{fancy}
\fancyhf{}
\fancyhead[L]{\small Loan Default Prediction}
\fancyhead[R]{\small GenAI Capstone Project}
\fancyfoot[C]{\thepage}

\lstset{
  basicstyle=\ttfamily\small,
  breaklines=true,
  frame=single,
  numbers=left,
  numberstyle=\tiny\color{gray},
  backgroundcolor=\color{gray!5}
}

\hypersetup{
  colorlinks=true,
  linkcolor=black,
  urlcolor=blue,
  citecolor=black
}

\onehalfspacing

\begin{document}

% ---- Title Page ----
\begin{titlepage}
  \centering
  \vspace*{2cm}

  {\LARGE\bfseries Loan Default Prediction System}\\[0.5cm]
  {\Large Using Machine Learning Classification Models}\\[2cm]

  {\large\textbf{Intro to GenAI --- Capstone Project Report}}\\[2.5cm]

  {\large\textbf{Team Members}}\\[0.5cm]
  \begin{tabular}{l l}
    \toprule
    \textbf{Name} & \textbf{Role} \\
    \midrule
    Dhanvin Vadlamudi   & Team Lead \\
    Meka Chaitanya Sai  & Team Member \\
    Killi Akshith Kumar & Team Member \\
    Akhil Nath Reddy    & Team Member \\
    \bottomrule
  \end{tabular}

  \vfill

  {\large March 2026}

\end{titlepage}

% ---- Abstract & Table of Contents ----
\newpage
\begin{abstract}
Loan defaults are a significant financial risk for lending institutions. In this project, we built an end-to-end machine learning system to predict whether a loan applicant is likely to default. We used a Kaggle dataset of 45,000 loan records with 13 features and trained three classification models: Logistic Regression (baseline), Decision Tree (intermediate), and XGBoost (advanced). After handling class imbalance and tuning hyperparameters, XGBoost achieved the best performance with 91.27\% accuracy, 0.82 F1-Score, and 0.98 ROC-AUC. The system was deployed as a Streamlit web application that allows users to input applicant profiles and receive real-time default probability predictions with risk classification.
\end{abstract}

\vspace{1cm}
\tableofcontents
\newpage

% ===================================================================
\section{Introduction}

Loan defaults cause massive losses for banks. Traditional manual checks are slow and miss patterns in data. Our goal is to build a machine learning system that:

\begin{itemize}[leftmargin=2cm]
  \item Predicts the \textbf{probability of default} based on applicant profiles.
  \item Classifies applicants into \textbf{Low}, \textbf{Medium}, and \textbf{High Risk} categories.
  \item Compares multiple ML models and identifies the best one.
  \item Deploys the system as a \textbf{Streamlit web application} for real-time predictions.
\end{itemize}

% ===================================================================
\section{Data Description}

We used the Loan Approval Classification dataset from Kaggle \cite{kaggle_dataset}, saved as \texttt{loan\_data.csv}, containing \textbf{45,000 records} with 14 columns.

\begin{table}[H]
  \centering
  \caption{Feature Overview}
  \begin{tabular}{l l p{7cm}}
    \toprule
    \textbf{Feature} & \textbf{Type} & \textbf{Description} \\
    \midrule
    person\_age             & Numerical   & Age of the applicant \\
    person\_gender          & Categorical & Gender (male/female) \\
    person\_education       & Categorical & Education level (High School, Associate, Bachelor, Master, Doctorate) \\
    person\_income          & Numerical   & Annual income (\$) \\
    person\_emp\_exp        & Numerical   & Employment experience (years) \\
    person\_home\_ownership & Categorical & Home ownership (RENT, OWN, MORTGAGE, OTHER) \\
    loan\_amnt              & Numerical   & Loan amount requested \\
    loan\_intent            & Categorical & Purpose of the loan (PERSONAL, EDUCATION, MEDICAL, VENTURE, HOMEIMPROVEMENT, DEBTCONSOLIDATION) \\
    loan\_int\_rate         & Numerical   & Interest rate (\%) \\
    loan\_percent\_income   & Numerical   & Loan amount as \% of income \\
    cb\_person\_cred\_hist\_length & Numerical & Credit history length (years) \\
    credit\_score           & Numerical   & Credit score (300--850) \\
    previous\_loan\_defaults\_on\_file & Categorical & Defaulted before (Yes/No) \\
    \midrule
    \textbf{loan\_status}   & \textbf{Target} & \textbf{1 = Default, 0 = No Default} \\
    \bottomrule
  \end{tabular}
\end{table}

The dataset is \textbf{imbalanced}: 35,000 records (78\%) are No Default and 10,000 records (22\%) are Default.

% ===================================================================
\section{Exploratory Data Analysis (EDA)}

We performed thorough Exploratory Data Analysis (EDA) in \texttt{preprocessing.ipynb} to understand the underlying distributions and relationships in the dataset prior to modelling. Our analysis included visualising data distributions with histograms and box plots, as well as examining feature interactions through a correlation matrix heatmap.

\textbf{Missing Values:} We identified missing values in numerical columns and opted to drop rows with nulls using \texttt{df.dropna()} to ensure data integrity without introducing bias from imputation.

\textbf{Outliers:} We observed significant rightward skewness in several financial features. We capped extreme values at the 99th percentile for \texttt{person\_age}, \texttt{person\_income}, and \texttt{person\_emp\_exp} using \texttt{df['column'].quantile(0.99)}, standardising the scale for our models.

\textbf{Key Findings from Correlation Analysis:}
\begin{itemize}[leftmargin=2cm]
  \item \texttt{loan\_percent\_income} had the highest positive correlation with default, indicating applicants borrowing a large portion of their income are highly risky.
  \item \texttt{previous\_loan\_defaults\_on\_file} was strongly correlated --- past defaulters have a significantly higher propensity to default again.
  \item \texttt{credit\_score} and \texttt{person\_income} had strong negative correlations --- predictably, higher values correspond to lower risk.
\end{itemize}

% ===================================================================
\section{Data Preprocessing}

Robust data preprocessing is essential for training accurate and generalisable models. All preprocessing steps are encapsulated in our \texttt{preprocessing.py} script:

\begin{enumerate}[leftmargin=2cm]
  \item \textbf{Cleaning}: Dropped nulls to maintain pristine data quality, capped outliers at the 99th percentile to prevent models from learning anomalous financial extremities, and binary-encoded gender (male=1, female=0) and previous defaults (Yes=1, No=0).
  \item \textbf{Ordinal Encoding}: Education mapped to a distinct 0--4 scale (High School < Associate < Bachelor < Master < Doctorate) to preserve the implicit hierarchy of qualifications.
  \item \textbf{One-Hot Encoding}: Nominal categorical variables such as \texttt{person\_home\_ownership} and \texttt{loan\_intent} were transformed into sparse matrices via \texttt{pd.get\_dummies(drop\_first=True)}, avoiding the dummy variable trap.
  \item \textbf{Train-Test Split}: We partitioned the data using an 80/20 split with \texttt{random\_state=42}, ensuring a robust test set for realistic evaluation metrics.
  \item \textbf{Feature Scaling}: \texttt{StandardScaler} was applied exclusively for Logistic Regression to converge faster, whereas tree-based models were trained on unscaled data since they rely on value splitting rather than distance metrics.
\end{enumerate}

% ===================================================================
\section{Methodology}

We trained three classification models, progressing from simple to advanced:

\subsection{Logistic Regression (Baseline)}

A linear model using the sigmoid function to output default probability \cite{sklearn}. The predicted probability is:

\begin{equation}
P(y=1 \mid x) = \sigma(w^T x + b) = \frac{1}{1 + e^{-(w^T x + b)}}
\end{equation}

Configured with \texttt{max\_iter=1000}, \texttt{class\_weight='balanced'}, and \texttt{StandardScaler}. Chosen as our baseline because it is fast, interpretable, and gives us a reference point.

\subsection{Decision Tree (Intermediate)}

A non-linear model that splits data at feature thresholds using Gini Impurity \cite{sklearn}:

\begin{equation}
Gini = 1 - \sum_{i=1}^{C} p_i^2
\end{equation}

where $p_i$ is the proportion of class $i$ in a node. Configured with \texttt{max\_depth=8}, \texttt{min\_samples\_leaf=10}, and \texttt{class\_weight='balanced'}. Chosen because it captures feature interactions that logistic regression cannot.

\subsection{XGBoost (Best Model)}

A gradient boosting algorithm that builds trees sequentially, each minimising the objective function \cite{xgboost}:

\begin{equation}
\mathcal{L} = \sum_{i=1}^{n} l(y_i, \hat{y}_i) + \sum_{k=1}^{K} \Omega(f_k)
\end{equation}

where $l$ is the log-loss function and $\Omega$ is the regularisation term that prevents overfitting. Key parameters:

\begin{itemize}[leftmargin=2cm]
  \item \texttt{n\_estimators=200}, \texttt{max\_depth=8}, \texttt{learning\_rate=0.05}
  \item \texttt{scale\_pos\_weight} --- auto-calculated ratio to handle class imbalance
  \item \texttt{subsample=0.8}, \texttt{colsample\_bytree=0.8} --- regularisation to prevent overfitting
\end{itemize}

Chosen because XGBoost is the best-performing algorithm for structured/tabular datasets.

% ===================================================================
\section{Results}

All models were evaluated on the \textbf{test set} (20\% unseen data) using the following metrics:

\begin{equation}
\text{Accuracy} = \frac{TP + TN}{TP + TN + FP + FN}
\end{equation}

\begin{equation}
\text{F1-Score} = 2 \cdot \frac{\text{Precision} \cdot \text{Recall}}{\text{Precision} + \text{Recall}}
\end{equation}

\textbf{ROC-AUC} measures the area under the Receiver Operating Characteristic curve --- a value of 1.0 means perfect classification.

\begin{table}[H]
  \centering
  \caption{Model Comparison on Test Set}
  \label{tab:results}
  \renewcommand{\arraystretch}{1.5}
  \begin{tabular}{l c c c}
    \toprule
    \rowcolor{blue!20} \textbf{Model} & \textbf{Accuracy (\%)} & \textbf{F1-Score} & \textbf{ROC-AUC} \\
    \midrule
    Logistic Regression & 85.25 & 0.7389 & 0.9523 \\
    \rowcolor{gray!10} Decision Tree       & 88.03 & 0.7714 & 0.9627 \\
    \rowcolor{green!20} \textbf{XGBoost}    & \textbf{91.27} & \textbf{0.8239} & \textbf{0.9763} \\
    \bottomrule
  \end{tabular}
\end{table}

\textbf{XGBoost wins on every metric.} The 6\% accuracy jump from Logistic Regression to XGBoost shows that boosting captures complex patterns a linear model cannot. The high ROC-AUC (0.98) means the model separates defaulters from non-defaulters almost perfectly.

\textbf{Top predictive features} across all models: \texttt{loan\_percent\_income}, \texttt{previous\_loan\_defaults\_on\_file}, \texttt{credit\_score}, \texttt{loan\_int\_rate}, and \texttt{person\_income}.

% ===================================================================
\section{Optimisation}

Steps taken to improve model performance:

\begin{enumerate}[leftmargin=2cm]
  \item \textbf{Class Imbalance Handling}: Used \texttt{class\_weight='balanced'} in Logistic Regression and Decision Tree, and \texttt{scale\_pos\_weight} in XGBoost. Without this, the model would just predict ``No Default'' for everyone and get 78\% accuracy --- which is useless.
  \item \textbf{Hyperparameter Selection}: We tried different values of \texttt{max\_depth} (5, 8, 10) and found 8 gave the best balance. For XGBoost, \texttt{learning\_rate=0.05} with 200 trees performed better than the default settings. We also added \texttt{min\_samples\_leaf=10} to prevent the Decision Tree from creating leaves with too few samples.
  \item \textbf{Outlier Removal}: Capping at the 99th percentile helped models focus on realistic data ranges instead of being skewed by extreme values.
  \item \textbf{Feature Scaling}: \texttt{StandardScaler} improved Logistic Regression convergence since it is sensitive to feature magnitudes.
\end{enumerate}

% ===================================================================
\section{Application --- Streamlit Dashboard}

We built a web application using Streamlit \cite{streamlit} with the following features:

\textbf{Hosted App:} \url{https://genaicapstone-m8ezlrzqqq8ctmmfgsh6ur.streamlit.app/}

\textbf{GitHub Repository:} \url{https://github.com/ChaitanyaSai-Meka/gen_ai_capstone}

\begin{enumerate}[leftmargin=2cm]
  \item \textbf{Model Selection}: Users can choose any of the three models from the sidebar.
  \item \textbf{Model Comparison}: Bar chart (accuracy) and ROC curves displayed side-by-side.
  \item \textbf{Real-Time Prediction}: Input applicant details via a form and get instant default probability with risk classification (Low/Medium/High).
  \item \textbf{Feature Importance}: After prediction, shows the top 5 factors affecting the result.
  \item \textbf{Metrics Dashboard}: Sidebar shows active model's accuracy, F1-Score, and AUC.
\end{enumerate}

\begin{figure}[H]
  \centering
  \includegraphics[width=\textwidth]{dashboard.png}
  \caption{Streamlit Dashboard --- Model Comparison with Accuracy Bar Chart and ROC Curves}
  \label{fig:dashboard}
\end{figure}

\begin{figure}[H]
  \centering
  \includegraphics[width=\textwidth]{prediction_form.png}
  \caption{Real-Time Prediction Form with Risk Classification Output}
  \label{fig:prediction}
\end{figure}

\textbf{Tech Stack}: Streamlit, scikit-learn, XGBoost, Pandas, NumPy, Matplotlib.

% ===================================================================
\section{Team Contribution}

\begin{table}[H]
  \centering
  \caption{Team Contribution Breakdown}
  \begin{tabular}{l p{9cm}}
    \toprule
    \textbf{Team Member} & \textbf{Contribution} \\
    \midrule
    \textbf{Dhanvin Vadlamudi} (Team Lead)
  & Led the overall project architecture, technical direction, and planning. Designed and implemented the XGBoost model, including systematic hyperparameter tuning and performance optimization. Developed the end-to-end preprocessing pipeline. Structured the repository, established coding standards, and coordinated team tasks to ensure timely delivery and smooth integration. \\
    \addlinespace
    \textbf{Meka Chaitanya Sai}
      & Built the Decision Tree model. Experimented with different hyperparameters (\texttt{max\_depth}, \texttt{min\_samples\_leaf}). Assisted with feature encoding. \\
    \addlinespace
    \textbf{Killi Akshith Kumar}
      & Built the Logistic Regression model --- training, evaluation, and tuning. Performed initial EDA and data exploration in notebooks. Helped with data cleaning logic. \\
    \addlinespace
    \textbf{Akhil Nath Reddy}
      & Wrote the project report. Built the Streamlit application (\texttt{app.py}). Integrated all models into the dashboard. Assisted with testing and debugging. \\
    \bottomrule
  \end{tabular}
\end{table}

% ===================================================================
\section{Conclusion}

We built a complete loan default prediction system from scratch. We trained three models (Logistic Regression, Decision Tree, XGBoost), compared their performance, and found that \textbf{XGBoost} gave the best results with 91.27\% accuracy and 0.98 ROC-AUC.

Key takeaways:
\begin{itemize}[leftmargin=2cm]
  \item Handling class imbalance is critical --- without it, models just predict the majority class.
  \item XGBoost outperforms simpler models on tabular data thanks to gradient boosting.
  \item Deploying as a Streamlit app makes the project practical and demonstrable.
  \item \texttt{loan\_percent\_income} was consistently the top predictor of default.
\end{itemize}

% ===================================================================
\begin{thebibliography}{9}

\bibitem{hosted_app}
Streamlit Hosted Application, ``Loan Default Prediction Dashboard.'' \\
\url{https://genaicapstone-m8ezlrzqqq8ctmmfgsh6ur.streamlit.app/}

\bibitem{kaggle_dataset}
Kaggle Dataset, ``Loan Approval Classification Data.'' \\
\url{https://www.kaggle.com/datasets/taweilo/loan-approval-classification-data}

\bibitem{github_repo}
GitHub Repository, ``GenAI Capstone Source Code.'' \\
\url{https://github.com/ChaitanyaSai-Meka/gen_ai_capstone}

\end{thebibliography}

% ===================================================================
\section*{Academic Integrity Declaration}

We, the above-named team members, hereby affirm that the core logic, model architecture, preprocessing pipeline, and Streamlit application code in this project are our own original work. No Generative AI tool was used to directly produce the core implementation. All use of AI tools was limited to research, understanding concepts, and debugging, in accordance with course guidelines.

\vspace{1cm}

\noindent
\textbf{Team Lead:} Dhanvin Vadlamudi \\
\textbf{Date:} March 2026

\end{document}
